\chapter{Complexity — Oblivious Setup}
%%% generated by the blueprint pipeline from TCSlib.Complexity.TuringMachine.Robustness.ObliviousSetup
\section{Overview}

This module analyses the initialization (setup) phase of the length-only oblivious
schedule machine attached to a Turing machine $W$ over the Boolean alphabet and a
padding parameter $a$.  The explicit configuration \lean{Complexity.setupCfg} isolates
the machine's three controller tapes --- the binary budget tape, the unary macrostep
counter, and the sweep guide --- and phase-by-phase run identities give the exact cost
of the input reset, the copy of the input onto the guide, the fixed-width
binary-to-unary budget conversion, and the two-sided guide-zone allocation.  They
compose into \lean{Complexity.setupCfg_initializes}, the complete initialization run
with its exact duration.

\section{Declarations}

\begin{definition}[Explicit setup configuration]
\lean{Complexity.setupCfg}
\label{Complexity.setupCfg}
\leanok
\uses{Complexity.OblPhase, Complexity.OblSymbol, Turing.Cfg, Turing.FinTM}
Let $W$ be a Turing machine over the Boolean alphabet with $k$ work tapes and let
$a \in \bbn$ be a padding parameter; the oblivious schedule machine built from $W$ and
$a$ has $k+3$ work tapes over the schedule alphabet.  Given an input word $x$ of
schedule symbols, a base configuration $c$ on $x$, a controller phase $q$ (possibly the
halting state), an input-head position $p \in \{0,\dots,|x|+1\}$, integers
$b, u, g \in \bbz$, and tape contents $\beta, \mu, \gamma \colon \bbz \to
\Sigma \cup \{\text{blank}\}$ (where $\Sigma$ is the schedule alphabet), the associated
setup configuration is the configuration with state $q$ and input head at $p$, whose
first $k$ work tapes and head positions agree with those of $c$, whose last three
(controller) tapes --- the binary budget, unary macrostep counter, and sweep guide ---
have contents $\beta, \mu, \gamma$ with heads at $b, u, g$ respectively, and whose
output agrees with that of $c$.
\end{definition}

\begin{definition}[Optional write on one tape]
\lean{Complexity.setupWrite}
\label{Complexity.setupWrite}
\leanok
Given a tape $t \colon \bbz \to A \cup \{\text{blank}\}$, a cell $z \in \bbz$, and an
optional write instruction, the resulting tape is $t$ itself when no write is
requested, and otherwise $t$ updated at cell $z$ with the requested (possibly blank)
symbol.
\end{definition}

\begin{lemma}[Action of one schedule step on a setup configuration]
\lean{Complexity.setupCfg_apply}
\label{Complexity.setupCfg_apply}
\leanok
\uses{Complexity.OblPhase, Complexity.OblSymbol, Complexity.finThree_cases, Complexity.oblAction, Complexity.setupCfg, Complexity.setupWrite, Turing.Action.apply, Turing.Cfg, Turing.FinTM, Turing.moveInputPos}
\difficulty{4}
Consider a setup configuration of the oblivious schedule machine built from a Boolean
machine $W$ with $k$ work tapes and padding parameter $a$, with input head at $p$ and
controller tapes $\beta, \mu, \gamma$ with heads at $b, u, g$.  Applying a schedule
action with next phase $q'$, input-head move $d \in \{-1,0,+1\}$, and one (optional
write, move) pair for each controller tape yields exactly the setup configuration with
phase $q'$, input head moved by $d$ (clamped to $\{0,\dots,|x|+1\}$), each controller
tape updated by its optional write at the old head position, and each controller head
shifted by its move; the first $k$ work tapes and the output are untouched.
\end{lemma}

\begin{lemma}[Input symbol read in a setup configuration]
\lean{Complexity.setupCfg_input}
\label{Complexity.setupCfg_input}
\leanok
\uses{Complexity.OblPhase, Complexity.OblSymbol, Complexity.setupCfg, Turing.Cfg, Turing.Cfg.inputSymbol, Turing.FinTM}
\difficulty{2}
In a setup configuration on input $x$ with input head at position
$p \in \{0,\dots,|x|+1\}$, the symbol read from the input tape is blank when $p = 0$,
and otherwise it is the $(p-1)$-st symbol of $x$ (in $0$-indexed order), which is blank
when $p - 1 \ge |x|$.
\end{lemma}

\begin{lemma}[Controller tape reads in a setup configuration]
\lean{Complexity.setupCfg_reads}
\label{Complexity.setupCfg_reads}
\leanok
\uses{Complexity.OblPhase, Complexity.OblSymbol, Complexity.finThree_cases, Complexity.setupCfg, Turing.Cfg, Turing.Cfg.workTapeSymbols, Turing.FinTM}
\difficulty{3}
In a setup configuration with controller tapes $\beta, \mu, \gamma$ and heads at
$b, u, g$, the symbols read from the three controller tapes are $\beta(b)$, $\mu(u)$,
and $\gamma(g)$ respectively.
\end{lemma}

\begin{lemma}[Exact cost of the reset scan]
\lean{Complexity.setupCfg_reset_scan}
\label{Complexity.setupCfg_reset_scan}
\leanok
\uses{Complexity.OblPhase, Complexity.OblSymbol, Complexity.oblAction, Complexity.obliviousSchedule, Complexity.setupCfg, Complexity.setupCfg_apply, Complexity.setupCfg_input, Complexity.setupWrite, Turing.Action.apply, Turing.Cfg, Turing.Cfg.inputSymbol, Turing.FinTM, Turing.MultiTapeTM.runFrom, Turing.MultiTapeTM.runFrom_succ_eq_step, Turing.moveInputPos}
\difficulty{5}
Suppose the oblivious schedule machine is in a setup configuration in the reset-scan
phase with input head at an interior position $j \le |x|$; the controller tapes and
heads are arbitrary.  Then after exactly $j$ steps the machine is in the same setup
configuration except that the input head is at the left boundary (position $0$): the
phase, all tapes, and all work heads are unchanged.
\end{lemma}

\begin{lemma}[Exact cost of the input reset]
\lean{Complexity.setupCfg_reset}
\label{Complexity.setupCfg_reset}
\leanok
\uses{Complexity.OblPhase, Complexity.OblSymbol, Complexity.copyGuide, Complexity.obliviousSchedule, Complexity.setupCfg, Complexity.setupCfg_apply, Complexity.setupCfg_input, Complexity.setupCfg_reset_scan, Complexity.setupWrite, Turing.Action.apply, Turing.Cfg, Turing.FinTM, Turing.MultiTapeTM.runFrom, Turing.MultiTapeTM.runFrom_add, Turing.MultiTapeTM.runFrom_succ_eq_step, Turing.MultiTapeTM.runFrom_zero, Turing.MultiTapeTM.step, Turing.moveInputPos}
\difficulty{5}
From a setup configuration in the reset-start phase with input head at an arbitrary
position $p \in \{0,\dots,|x|+1\}$, running exactly $\max(p,1)+1$ steps of the
oblivious schedule machine leads to the setup configuration in the copy-left phase with
input head at position $1$; all controller tapes and heads are unchanged.  (The reset
begins with a mandatory left move, so even the right blank boundary and the empty input
reach the left boundary before the final right move.)
\end{lemma}

\begin{definition}[Partially copied guide tape]
\lean{Complexity.copyGuide}
\label{Complexity.copyGuide}
\leanok
\uses{Complexity.OblSymbol}
For $j \in \bbn$, the partial guide after copying $j$ input cells is the tape whose
cell $-1$ carries an interior mark (the virtual left boundary), whose cells $z$ with
$0 \le z < j$ carry the origin mark at $z = 0$ and interior marks otherwise, and whose
remaining cells are blank.
\end{definition}

\begin{lemma}[Extending the partial guide by one cell]
\lean{Complexity.copyGuide_next}
\label{Complexity.copyGuide_next}
\leanok
\uses{Complexity.copyGuide}
\difficulty{3}
Updating the partial guide with $j$ copied cells at cell $j$ --- writing the origin
mark when $j = 0$ and an interior mark otherwise --- yields the partial guide with
$j+1$ copied cells.
\end{lemma}

\begin{lemma}[Origin uniqueness in the partial guide]
\lean{Complexity.copyGuide_origin}
\label{Complexity.copyGuide_origin}
\leanok
\uses{Complexity.copyGuide}
\difficulty{2}
For $j > 0$, a cell of the partial guide with $j$ copied cells carries the origin mark
if and only if it is cell $0$.
\end{lemma}

\begin{definition}[Copy loop phase]
\lean{Complexity.copyPhase}
\label{Complexity.copyPhase}
\leanok
\uses{Complexity.OblPhase}
The controller phase of the copy loop after $j$ copied symbols: the first-copy phase
when $j = 0$ (the origin mark has not yet been written) and the further-copy phase
otherwise.
\end{definition}

\begin{lemma}[The two steps opening the copy stage]
\lean{Complexity.setupCfg_copy_left}
\label{Complexity.setupCfg_copy_left}
\leanok
\uses{Complexity.OblPhase, Complexity.OblSymbol, Complexity.copyGuide, Complexity.copyPhase, Complexity.obliviousSchedule, Complexity.setupCfg, Complexity.setupCfg_apply, Complexity.setupWrite, Turing.Action.apply, Turing.Cfg, Turing.FinTM, Turing.MultiTapeTM.runFrom, Turing.MultiTapeTM.runFrom_succ_eq_step, Turing.MultiTapeTM.runFrom_zero, Turing.MultiTapeTM.step, Turing.moveInputPos_zero}
\difficulty{4}
From a setup configuration in the copy-left phase with guide head at $0$ and entirely
blank guide tape, exactly two steps of the oblivious schedule machine write the virtual
left boundary (an interior mark at guide cell $-1$) and reach the first-copy phase with
guide head back at $0$ over the partial guide with $0$ copied cells; the input head,
the budget and unary tapes, and their heads are unchanged.  This holds for every
input-head position, including the empty input.
\end{lemma}

\begin{lemma}[Prefix of the copy scan]
\lean{Complexity.setupCfg_copy_prefix}
\label{Complexity.setupCfg_copy_prefix}
\leanok
\uses{Complexity.OblPhase, Complexity.OblSymbol, Complexity.copyGuide, Complexity.copyGuide_next, Complexity.copyPhase, Complexity.oblAction, Complexity.obliviousSchedule, Complexity.setupCfg, Complexity.setupCfg_apply, Complexity.setupCfg_input, Complexity.setupWrite, Turing.Action.apply, Turing.Cfg, Turing.Cfg.inputSymbol, Turing.FinTM, Turing.MultiTapeTM.runFrom, Turing.MultiTapeTM.runFrom_succ_eq_step', Turing.moveInputPos}
\difficulty{5}
Let $j \le |x|$.  From a setup configuration in the first-copy phase with input head at
position $1$ and guide head at $0$ over the partial guide with $0$ copied cells,
exactly $j$ steps lead to the copy phase for $j$ copied symbols with input head at
$j+1$, guide head at $j$, and guide tape the partial guide with $j$ copied cells; the
budget and unary tapes and heads are unchanged.  In particular the input and guide
heads make exactly the same number of right moves, and no data value influences a
movement.
\end{lemma}

\begin{lemma}[The final blank read of the copy scan]
\lean{Complexity.setupCfg_copy_end}
\label{Complexity.setupCfg_copy_end}
\leanok
\uses{Complexity.OblPhase, Complexity.OblSymbol, Complexity.copyGuide, Complexity.copyGuide_next, Complexity.copyPhase, Complexity.oblAction, Complexity.obliviousSchedule, Complexity.setupCfg, Complexity.setupCfg_apply, Complexity.setupCfg_input, Complexity.setupWrite, Turing.Action.apply, Turing.Cfg, Turing.Cfg.inputSymbol, Turing.FinTM, Turing.MultiTapeTM.step, Turing.moveInputPos_zero}
\difficulty{4}
From a setup configuration in the copy phase for $|x|$ copied symbols, with input head
parked at the right boundary $|x|+1$ (reading blank) and guide head at $|x|$, a single
step enters the copy-return phase, leaves both heads in place, and closes the copied
zone by writing one more mark, producing the partial guide with $|x|+1$ cells.
\end{lemma}

\begin{lemma}[Return across the copied guide]
\lean{Complexity.setupCfg_copy_return}
\label{Complexity.setupCfg_copy_return}
\leanok
\uses{Complexity.OblPhase, Complexity.OblSymbol, Complexity.copyGuide, Complexity.copyGuide_origin, Complexity.obliviousSchedule, Complexity.setupCfg, Complexity.setupCfg_apply, Complexity.setupWrite, Complexity.unaryTape, Turing.Action.apply, Turing.Cfg, Turing.Cfg.workTapeSymbols, Turing.FinTM, Turing.MultiTapeTM.runFrom, Turing.MultiTapeTM.runFrom_succ_eq_step, Turing.MultiTapeTM.runFrom_zero, Turing.moveInputPos_zero}
\difficulty{5}
For all $n, j \in \bbn$: from a setup configuration in the copy-return phase with guide
head at $j$ over the partial guide with $n+1$ copied cells, exactly $j+1$ steps reach
the budget-start phase with guide head at the origin (cell $0$); all tapes and all
other heads are unchanged.
\end{lemma}

\begin{definition}[Unary macrostep counter tape]
\lean{Complexity.unaryTape}
\label{Complexity.unaryTape}
\leanok
\uses{Complexity.OblSymbol}
For $N \in \bbn$, the unary tape with $N$ markers has the origin sentinel at cell $0$,
unit markers on cells $1,\dots,N$, and blanks elsewhere.
\end{definition}

\begin{lemma}[Appending one unary marker]
\lean{Complexity.unaryTape_next}
\label{Complexity.unaryTape_next}
\leanok
\uses{Complexity.unaryTape}
\difficulty{3}
Updating the unary tape with $N$ markers at cell $N+1$ with a unit marker yields the
unary tape with $N+1$ markers.
\end{lemma}

\begin{lemma}[Origin uniqueness on the captured clock tape]
\lean{Complexity.clockTape_origin}
\label{Complexity.clockTape_origin}
\leanok
\uses{Complexity.clockTape, Turing.FinTM.bufferTape}
\difficulty{3}
Let $w$ be a bit word and consider its captured clock tape: the tape with the origin
sentinel at cell $0$, the $i$-th bit of $w$ at cell $i+1$ for $0 \le i < |w|$, and
blanks elsewhere.  A cell of this tape carries the origin mark if and only if it is
cell $0$.
\end{lemma}

\begin{lemma}[Blank cell after the captured word]
\lean{Complexity.clockTape_end}
\label{Complexity.clockTape_end}
\leanok
\uses{Complexity.clockTape}
\difficulty{1}
On the captured clock tape of a bit word $w$ (origin sentinel at cell $0$, the bits of
$w$ on cells $1,\dots,|w|$), cell $|w|+1$ is blank.
\end{lemma}

\begin{lemma}[Rewind of the budget tape]
\lean{Complexity.setupCfg_budget_back}
\label{Complexity.setupCfg_budget_back}
\leanok
\uses{Complexity.OblPhase, Complexity.OblSymbol, Complexity.clockTape, Complexity.clockTape_origin, Complexity.obliviousSchedule, Complexity.setupCfg, Complexity.setupCfg_apply, Complexity.setupWrite, Turing.Action.apply, Turing.Cfg, Turing.Cfg.workTapeSymbols, Turing.FinTM, Turing.MultiTapeTM.runFrom, Turing.MultiTapeTM.runFrom_succ_eq_step, Turing.MultiTapeTM.runFrom_zero, Turing.MultiTapeTM.step}
\difficulty{5}
For every $j \in \bbn$: from a setup configuration in the budget-back phase with budget
head at cell $j$ over the captured clock tape of a bit word $w$, exactly $j+1$ steps
reach the sentinel at the origin and move right once, entering the first append phase
with budget head at cell $1$; all tapes and all other heads are unchanged.
\end{lemma}

\begin{lemma}[Exact cost of a full budget rewind]
\lean{Complexity.setupCfg_budget_start}
\label{Complexity.setupCfg_budget_start}
\leanok
\uses{Complexity.OblPhase, Complexity.OblSymbol, Complexity.clockTape, Complexity.obliviousSchedule, Complexity.setupCfg, Complexity.setupCfg_apply, Complexity.setupCfg_budget_back, Complexity.setupWrite, Turing.Action.apply, Turing.Cfg, Turing.FinTM, Turing.MultiTapeTM.runFrom, Turing.MultiTapeTM.runFrom_succ_eq_step, Turing.MultiTapeTM.step}
\difficulty{4}
From a setup configuration in the budget-start phase with budget head at cell $|w|+1$
(one cell beyond the captured clock tape of the bit word $w$), exactly $|w|+2$ steps
reach the first append phase with budget head at cell $1$; the cost depends only on the
width of $w$, not on its bits, and all tapes and all other heads are unchanged.
\end{lemma}

\begin{lemma}[Prefix of the unary append block]
\lean{Complexity.setupCfg_append_prefix}
\label{Complexity.setupCfg_append_prefix}
\leanok
\uses{Complexity.OblPhase, Complexity.OblSymbol, Complexity.obliviousSchedule, Complexity.setupCfg, Complexity.setupCfg_apply, Complexity.setupWrite, Complexity.unaryTape, Complexity.unaryTape_next, Turing.Action.apply, Turing.Cfg, Turing.FinTM, Turing.MultiTapeTM.runFrom, Turing.MultiTapeTM.runFrom_succ_eq_step', Turing.moveInputPos_zero}
\difficulty{5}
Let $j \le a+1$, where $a$ is the padding parameter.  From a setup configuration in
append phase $0$ with unary head at cell $N+1$ over the unary tape with $N$ markers,
exactly $j$ steps reach append phase $j$ with unary head at cell $N+j+1$ over the unary
tape with $N+j$ markers; each step writes exactly one unit marker, and all other tapes
and heads are unchanged.
\end{lemma}

\begin{lemma}[The full unary append block]
\lean{Complexity.setupCfg_append}
\label{Complexity.setupCfg_append}
\leanok
\uses{Complexity.OblPhase, Complexity.OblSymbol, Complexity.obliviousSchedule, Complexity.setupCfg, Complexity.setupCfg_append_prefix, Complexity.setupCfg_apply, Complexity.setupWrite, Complexity.unaryTape, Turing.Action.apply, Turing.Cfg, Turing.FinTM, Turing.MultiTapeTM.runFrom, Turing.MultiTapeTM.runFrom_succ_eq_step'}
\difficulty{3}
From a setup configuration in append phase $0$ with unary head at cell $N+1$ over the
unary tape with $N$ markers, exactly $a+2$ steps append exactly $a+1$ unit markers and
enter the borrow phase with incoming borrow set: the unary head is then at cell
$N+(a+1)+1$ over the unary tape with $N+a+1$ markers, and all other tapes and heads are
unchanged.
\end{lemma}

\begin{lemma}[Shifting a tape zipper past a fixed word]
\lean{Complexity.sweepTape_shift}
\label{Complexity.sweepTape_shift}
\leanok
\uses{Turing.FinTM.sweepTape, Turing.FinTM.sweepTape_read, Turing.FinTM.sweepTape_right}
\difficulty{5}
Present a tape as a zipper with frontier $z \in \bbz$, a list $l$ of the (possibly
blank) cells to the left of $z$ listed nearest-first, and a list of the cells at and to
the right of $z$.  If the right list decomposes as $w \mathbin{+\!\!+} r$, then the
same tape is presented by the zipper with frontier $z + |w|$, left list
$\mathrm{reverse}(w) \mathbin{+\!\!+} l$, and right list $r$.
\end{lemma}

\begin{lemma}[The captured clock tape as a zipper]
\lean{Complexity.clockTape_zipper}
\label{Complexity.clockTape_zipper}
\leanok
\uses{Complexity.OblSymbol, Complexity.clockTape, Turing.FinTM.bufferTape, Turing.FinTM.sweepTape}
\difficulty{5}
For every bit word $w$, the captured clock tape of $w$ (origin sentinel at cell $0$,
the bits of $w$ on cells $1,\dots,|w|$, blanks elsewhere) equals the zipper tape with
frontier at cell $1$, left list consisting of the single origin sentinel, and right
list the bits of $w$, each encoded as a budget-bit symbol.
\end{lemma}

\begin{lemma}[The end zipper of a captured word]
\lean{Complexity.clockTape_zipper_end}
\label{Complexity.clockTape_zipper_end}
\leanok
\uses{Complexity.OblSymbol, Complexity.clockTape, Complexity.clockTape_zipper, Complexity.sweepTape_shift, Turing.FinTM.sweepTape}
\difficulty{2}
For every bit word $w$, the zipper tape with frontier $1+|w|$, empty right list, and
left list consisting of the reversed encoded bits of $w$ followed by the origin
sentinel equals the captured clock tape of $w$.
\end{lemma}

\begin{lemma}[Lane focusing of a setup configuration]
\lean{Complexity.laneCfg_setup}
\label{Complexity.laneCfg_setup}
\leanok
\uses{Complexity.OblPhase, Complexity.OblSymbol, Complexity.finThree_cases, Complexity.laneCfg, Complexity.setupCfg, Turing.Cfg, Turing.FinTM, Turing.FinTM.sweepTape}
\difficulty{4}
Take a setup configuration and focus its budget lane: replace the state by $q'$, the
budget tape by the zipper tape with frontier $z$, left list $l$, and right list $r$,
and move the budget head to $z$, keeping every other tape and head.  The result is
again a setup configuration, with the same base configuration, input head, unary and
guide data, state $q'$, budget head $z$, and budget tape given by that zipper.
\end{lemma}

\begin{lemma}[The full-width borrow pass]
\lean{Complexity.setupCfg_borrow}
\label{Complexity.setupCfg_borrow}
\leanok
\uses{Complexity.OblPhase, Complexity.OblSymbol, Complexity.budgetBorrow, Complexity.budgetBorrow_length, Complexity.clockTape, Complexity.clockTape_zipper, Complexity.clockTape_zipper_end, Complexity.laneCfg_setup, Complexity.obliviousSchedule, Complexity.obliviousSchedule_borrow_run, Complexity.setupCfg, Turing.Cfg, Turing.FinTM, Turing.MultiTapeTM.runFrom}
\difficulty{4}
Let $w$ be a bit word and $c$ a Boolean incoming borrow, and recall the fixed-width
borrow pass: it processes all of $w$, subtracting the incoming borrow from its
little-endian value while preserving its length, and returns an outgoing borrow flag
together with the rewritten word.  From a setup configuration in the borrow phase with
flag $c$ and budget head at cell $1$ over the captured clock tape of $w$, exactly $|w|$
steps reach the borrow phase whose flag is the outgoing borrow, with budget head at
$1+|w|$ over the captured clock tape of the rewritten word; all other tapes and heads,
and the parked input head, remain fixed.
\end{lemma}

\begin{lemma}[Branch at the end of a borrow pass]
\lean{Complexity.setupCfg_borrow_end}
\label{Complexity.setupCfg_borrow_end}
\leanok
\uses{Complexity.OblPhase, Complexity.OblSymbol, Complexity.clockTape, Complexity.clockTape_end, Complexity.obliviousSchedule, Complexity.setupCfg, Complexity.setupCfg_apply, Complexity.setupWrite, Turing.Action.apply, Turing.Cfg, Turing.Cfg.workTapeSymbols, Turing.FinTM, Turing.MultiTapeTM.step}
\difficulty{4}
From a setup configuration in the borrow phase with flag $c$ and budget head at cell
$|w|+1$ (the blank cell after the captured clock tape of the bit word $w$), one step
keeps every tape and head fixed and moves to the unary-start phase when $c$ is true
(the budget underflowed, i.e.\ is exhausted) and back to the budget-start phase
otherwise.
\end{lemma}

\begin{lemma}[Exact cost of one budget round]
\lean{Complexity.setupCfg_budget_round}
\label{Complexity.setupCfg_budget_round}
\leanok
\uses{Complexity.OblPhase, Complexity.OblSymbol, Complexity.budgetBorrow, Complexity.budgetBorrow_length, Complexity.clockTape, Complexity.obliviousSchedule, Complexity.setupCfg, Complexity.setupCfg_append, Complexity.setupCfg_borrow, Complexity.setupCfg_borrow_end, Complexity.setupCfg_budget_start, Complexity.unaryTape, Turing.Cfg, Turing.FinTM, Turing.MultiTapeTM.runFrom, Turing.MultiTapeTM.runFrom_add, Turing.MultiTapeTM.runFrom_succ_eq_step'}
\difficulty{4}
One complete round of the budget controller --- full rewind, append block of $a+1$
unary markers, fixed-width borrow pass, and end test --- costs exactly
$2|w| + (a+1) + 4$ steps, independently of the bits of the budget word $w$.  Starting
in the budget-start phase with budget head at $|w|+1$ over the captured clock tape of
$w$ and unary head at $N+1$ over the unary tape with $N$ markers, the round ends with
the captured clock tape of the borrowed word (incoming borrow set), the unary tape with
$N + a + 1$ markers, and phase unary-start or budget-start according to the borrow's
underflow flag.
\end{lemma}

\begin{lemma}[The complete binary-to-unary conversion]
\lean{Complexity.setupCfg_budget_all}
\label{Complexity.setupCfg_budget_all}
\leanok
\uses{Complexity.OblPhase, Complexity.OblSymbol, Complexity.budgetBorrow, Complexity.budgetBorrow_length, Complexity.budgetBorrow_underflow, Complexity.budgetBorrow_value, Complexity.budgetValue, Complexity.clockTape, Complexity.obliviousSchedule, Complexity.setupCfg, Complexity.setupCfg_budget_round, Complexity.unaryTape, Turing.Cfg, Turing.FinTM, Turing.MultiTapeTM.runFrom, Turing.MultiTapeTM.runFrom_add}
\difficulty{5}
Let $w$ be a bit word of little-endian value $v$ and let $N \in \bbn$.  From a setup
configuration in the budget-start phase with budget head at $|w|+1$ over the captured
clock tape of $w$ and unary head at $N+1$ over the unary tape with $N$ markers, exactly
$(v+1)\,\bigl(2|w| + (a+1) + 4\bigr)$ steps reach the unary-start phase with unary head
at $N + (a+1)(v+1) + 1$ over the unary tape with $N + (a+1)(v+1)$ markers, the budget
tape being the captured clock tape of some bit word $w'$ with $|w'| = |w|$ (the word
keeps its full width as its high bits become zero); the guide tape and head and the
input head are unchanged.
\end{lemma}

\begin{lemma}[Origin uniqueness on the unary tape]
\lean{Complexity.unaryTape_origin}
\label{Complexity.unaryTape_origin}
\leanok
\uses{Complexity.unaryTape}
\difficulty{2}
A cell of the unary tape with $N$ markers carries the origin mark if and only if it is
cell $0$.
\end{lemma}

\begin{lemma}[Unit markers of the unary tape]
\lean{Complexity.unaryTape_unit}
\label{Complexity.unaryTape_unit}
\leanok
\uses{Complexity.unaryTape}
\difficulty{1}
For every $j < N$, cell $j+1$ of the unary tape with $N$ markers carries a unit marker.
\end{lemma}

\begin{lemma}[Blank cell terminating the unary segment]
\lean{Complexity.unaryTape_end}
\label{Complexity.unaryTape_end}
\leanok
\uses{Complexity.unaryTape}
\difficulty{1}
Cell $N+1$ of the unary tape with $N$ markers is blank.
\end{lemma}

\begin{definition}[Common phase of the two unary rewinds]
\lean{Complexity.unaryRewindPhase}
\label{Complexity.unaryRewindPhase}
\leanok
\uses{Complexity.OblPhase}
The controller phase used to rewind the unary counter: the unary-reset phase when the
Boolean flag marks the final rewind, and the unary-back phase otherwise.  Both rewinds
follow the same path and differ only in their successor phase.
\end{definition}

\begin{lemma}[Exact cost of a unary rewind]
\lean{Complexity.setupCfg_unary_rewind}
\label{Complexity.setupCfg_unary_rewind}
\leanok
\uses{Complexity.OblPhase, Complexity.OblSymbol, Complexity.guideMark, Complexity.obliviousSchedule, Complexity.setupCfg, Complexity.setupCfg_apply, Complexity.setupWrite, Complexity.unaryRewindPhase, Complexity.unaryTape, Complexity.unaryTape_origin, Turing.Action.apply, Turing.Cfg, Turing.Cfg.workTapeSymbols, Turing.FinTM, Turing.MultiTapeTM.runFrom, Turing.MultiTapeTM.runFrom_succ_eq_step, Turing.MultiTapeTM.runFrom_zero, Turing.MultiTapeTM.step}
\difficulty{5}
For every $j \in \bbn$ and either rewind variant: from a setup configuration in the
corresponding unary rewind phase with unary head at cell $j$ over the unary tape with
$N$ markers, exactly $j+1$ steps reach the successor phase --- the start-center phase
for the final rewind, the first rightward layout phase otherwise --- with unary head at
cell $1$; no tape is written and all other heads are unchanged.
\end{lemma}

\begin{lemma}[Opening rewind of the allocation stage]
\lean{Complexity.setupCfg_unary_start}
\label{Complexity.setupCfg_unary_start}
\leanok
\uses{Complexity.OblPhase, Complexity.OblSymbol, Complexity.guideMark, Complexity.obliviousSchedule, Complexity.setupCfg, Complexity.setupCfg_apply, Complexity.setupCfg_unary_rewind, Complexity.setupWrite, Complexity.unaryRewindPhase, Complexity.unaryTape, Turing.Action.apply, Turing.Cfg, Turing.FinTM, Turing.MultiTapeTM.runFrom, Turing.MultiTapeTM.runFrom_succ_eq_step, Turing.MultiTapeTM.step}
\difficulty{4}
From a setup configuration in the unary-start phase with unary head at cell $N+1$ over
the unary tape with $N$ markers, exactly $N+2$ steps rewind the entire newly
constructed unary segment and reach the first rightward layout phase with unary head at
cell $1$; all tapes and all other heads are unchanged.
\end{lemma}

\begin{definition}[Interior guide mark]
\lean{Complexity.guideMark}
\label{Complexity.guideMark}
\leanok
\uses{Complexity.OblSymbol}
The symbol written to guide cell $z$ during allocation: the origin mark when $z = 0$
and an interior mark otherwise.
\end{definition}

\begin{definition}[Partially allocated guide tape]
\lean{Complexity.guideFill}
\label{Complexity.guideFill}
\leanok
\uses{Complexity.OblSymbol, Complexity.guideMark}
Given a direction (rightward or leftward), a guide tape $\gamma$, and $j \in \bbn$, the
$j$-cell filled version of $\gamma$ is the tape agreeing with $\gamma$ except that each
cell $z$ with $0 \le z < j$ (rightward) respectively $0 \le -z < j$ (leftward) is
overwritten with the guide mark of $z$: the origin mark at $0$ and interior marks
elsewhere, independently of the previous contents.
\end{definition}

\begin{lemma}[Empty allocation prefix]
\lean{Complexity.guideFill_zero}
\label{Complexity.guideFill_zero}
\leanok
\uses{Complexity.OblSymbol, Complexity.guideFill}
\difficulty{2}
Filling zero cells of a guide tape, in either direction, leaves the tape unchanged.
\end{lemma}

\begin{lemma}[Origin preservation under allocation]
\lean{Complexity.guideFill_origin}
\label{Complexity.guideFill_origin}
\leanok
\uses{Complexity.OblSymbol, Complexity.guideFill, Complexity.guideMark}
\difficulty{2}
If a guide tape carries the origin mark exactly at cell $0$, then for every $j$ and
either direction, its $j$-cell filled version also carries the origin mark exactly at
cell $0$.
\end{lemma}

\begin{lemma}[One more allocated cell]
\lean{Complexity.guideFill_next}
\label{Complexity.guideFill_next}
\leanok
\uses{Complexity.OblSymbol, Complexity.guideFill, Complexity.guideMark}
\difficulty{4}
Updating the $j$-cell filled version of a guide tape at cell $j$ (rightward)
respectively $-j$ (leftward) with the corresponding guide mark yields the $(j+1)$-cell
filled version.
\end{lemma}

\begin{definition}[Phases of the outward allocation passes]
\lean{Complexity.layoutPhase}
\label{Complexity.layoutPhase}
\leanok
\uses{Complexity.OblPhase}
The finite controller phase of an outward allocation pass, indexed by
$i \in \{0,1,2\}$: the rightward layout phases when the direction flag is set and the
leftward layout phases otherwise.  The two passes are symmetric.
\end{definition}

\begin{definition}[Allocation head movement]
\lean{Complexity.layoutMove}
\label{Complexity.layoutMove}
\leanok
The guide-head movement of an allocation pass: one cell to the right for the rightward
pass and one cell to the left for the leftward pass, independently of the cell
contents.
\end{definition}

\begin{lemma}[One allocation step]
\lean{Complexity.setupCfg_layout_step}
\label{Complexity.setupCfg_layout_step}
\leanok
\uses{Complexity.OblPhase, Complexity.OblSymbol, Complexity.guideMark, Complexity.layoutMove, Complexity.layoutPhase, Complexity.nextThird, Complexity.obliviousSchedule, Complexity.setupCfg, Complexity.setupCfg_apply, Complexity.setupWrite, Turing.Action.apply, Turing.Cfg, Turing.Cfg.workTapeSymbols, Turing.FinTM, Turing.MultiTapeTM.step, Turing.moveInputPos_zero}
\difficulty{4}
Consider a setup configuration in layout phase $i \in \{0,1,2\}$ of either direction,
and suppose the unary tape reads a unit marker at the unary head and the guide tape
reads the origin mark at the guide head exactly when that head is at cell $0$.  Then
one step moves to layout phase $i+1 \bmod 3$, writes the guide mark at the guide head,
moves the guide head one cell in the chosen direction, and advances the unary head by
one exactly when $i = 2$; everything else is unchanged.  Thus the controller spends
three guide cells per unary marker.
\end{lemma}

\begin{lemma}[Allocation step on a filled guide]
\lean{Complexity.setupCfg_layout_fill}
\label{Complexity.setupCfg_layout_fill}
\leanok
\uses{Complexity.OblPhase, Complexity.OblSymbol, Complexity.guideFill, Complexity.guideFill_next, Complexity.guideFill_origin, Complexity.layoutMove, Complexity.layoutPhase, Complexity.nextThird, Complexity.obliviousSchedule, Complexity.setupCfg, Complexity.setupCfg_layout_step, Turing.Cfg, Turing.FinTM, Turing.MultiTapeTM.step}
\difficulty{3}
Let $\gamma$ be a guide tape carrying the origin mark exactly at cell $0$, and suppose
the unary tape reads a unit marker at the unary head.  Then one step from layout phase
$i$ (of either direction) with guide head at $\pm j$ over the $j$-cell filled version
of $\gamma$ leads to layout phase $i+1 \bmod 3$ with guide head at $\pm(j+1)$ over the
$(j+1)$-cell filled version, incrementing the unary head exactly when $i = 2$.
\end{lemma}

\begin{lemma}[Cyclic phase as a remainder]
\lean{Complexity.nextThird_mod}
\label{Complexity.nextThird_mod}
\leanok
\uses{Complexity.nextThird}
\difficulty{3}
For every $j \in \bbn$, the cyclic successor on $\{0,1,2\}$ of $j \bmod 3$ is
$(j+1) \bmod 3$.  Hence the finite allocation phase is the remainder modulo three of
the number of guide cells already written.
\end{lemma}

\begin{lemma}[Prefix of an outward allocation pass]
\lean{Complexity.setupCfg_layout_prefix}
\label{Complexity.setupCfg_layout_prefix}
\leanok
\uses{Complexity.OblPhase, Complexity.OblSymbol, Complexity.guideFill, Complexity.guideFill_zero, Complexity.layoutPhase, Complexity.nextThird_mod, Complexity.obliviousSchedule, Complexity.setupCfg, Complexity.setupCfg_layout_fill, Complexity.unaryTape, Complexity.unaryTape_unit, Turing.Cfg, Turing.FinTM, Turing.MultiTapeTM.runFrom, Turing.MultiTapeTM.runFrom_succ_eq_step'}
\difficulty{5}
Let $\gamma$ be a guide tape carrying the origin mark exactly at cell $0$, and let
$j \le 3B$.  From layout phase $0$ (of either direction) with unary head at cell $1$
over the unary tape with $B$ markers and guide head at cell $0$ over $\gamma$, exactly
$j$ steps lead to layout phase $j \bmod 3$ with unary head at
$\lfloor j/3 \rfloor + 1$, guide head at $\pm j$, and guide tape the $j$-cell filled
version of $\gamma$: the elapsed time equals the number of written guide cells, and the
unary head moves once every three cells.
\end{lemma}

\begin{lemma}[End of an outward allocation pass]
\lean{Complexity.setupCfg_layout_end}
\label{Complexity.setupCfg_layout_end}
\leanok
\uses{Complexity.OblPhase, Complexity.OblSymbol, Complexity.layoutMove, Complexity.layoutPhase, Complexity.obliviousSchedule, Complexity.rightGuide, Complexity.setupCfg, Complexity.setupCfg_apply, Complexity.setupWrite, Complexity.unaryTape, Complexity.unaryTape_end, Turing.Action.apply, Turing.Cfg, Turing.Cfg.workTapeSymbols, Turing.FinTM, Turing.MultiTapeTM.step}
\difficulty{4}
From layout phase $0$ (of either direction) with unary head at cell $B+1$ --- the
blank cell terminating the unary tape with $B$ markers --- and guide head at cell $g$,
one step writes an interior mark at guide cell $g$, moves the guide head one further
cell outward, keeps the unary head in place, and enters the right-edge or left-edge
phase according to the direction.
\end{lemma}

\begin{definition}[Right half of the completed guide]
\lean{Complexity.rightGuide}
\label{Complexity.rightGuide}
\leanok
\uses{Complexity.OblSymbol, Complexity.copyGuide}
For $R \in \bbn$, the right guide of radius $R$ is the tape obtained from the partial
guide with $R+1$ copied cells by writing the right outer boundary marker at cell $R+1$.
\end{definition}

\begin{lemma}[Rightward filling extends the copied guide]
\lean{Complexity.guideFill_right_copy}
\label{Complexity.guideFill_right_copy}
\leanok
\uses{Complexity.copyGuide, Complexity.guideFill, Complexity.guideMark}
\difficulty{2}
If $n + 1 \le R$, then filling $R$ cells rightward starting from the partial guide with
$n+1$ copied cells yields exactly the partial guide with $R$ cells: only the positive
prefix is overwritten.
\end{lemma}

\begin{lemma}[Origin uniqueness on the right guide]
\lean{Complexity.rightGuide_origin}
\label{Complexity.rightGuide_origin}
\leanok
\uses{Complexity.copyGuide_origin, Complexity.rightGuide}
\difficulty{3}
A cell of the right guide of radius $R$ carries the origin mark if and only if it is
cell $0$; the right-boundary write cannot alter the origin.
\end{lemma}

\begin{lemma}[Completion of the guide by the left pass]
\lean{Complexity.guideFill_left_finish}
\label{Complexity.guideFill_left_finish}
\leanok
\uses{Complexity.copyGuide, Complexity.guideFill, Complexity.guideMark, Complexity.guideTape, Complexity.rightGuide}
\difficulty{2}
Let $R > 0$.  Starting from the right guide of radius $R$, filling $R$ cells leftward
and then writing an interior mark at cell $-R$ and the left outer boundary marker at
cell $-R-1$ produces the completed guide of radius $R$: origin mark at $0$, interior
marks at all $z$ with $0 < |z| \le R$, boundary markers at $\pm(R+1)$, and blanks
elsewhere.
\end{lemma}

\begin{lemma}[The right boundary transition]
\lean{Complexity.setupCfg_right_edge}
\label{Complexity.setupCfg_right_edge}
\leanok
\uses{Complexity.OblPhase, Complexity.OblSymbol, Complexity.copyGuide, Complexity.obliviousSchedule, Complexity.rightGuide, Complexity.setupCfg, Complexity.setupCfg_apply, Complexity.setupWrite, Complexity.unaryTape, Turing.Action.apply, Turing.Cfg, Turing.FinTM, Turing.MultiTapeTM.step}
\difficulty{3}
From a setup configuration in the right-edge phase with unary head at cell $B+1$ over
the unary tape with $B$ markers and guide head at cell $R+1$ over the partial guide
with $R+1$ copied cells, one step writes the right outer boundary marker at guide cell
$R+1$, moves the unary and guide heads one cell left (to $B$ and $R$), and enters the
first layout-return phase; the resulting guide tape is the right guide of radius $R$.
\end{lemma}

\begin{lemma}[One step of the read-only return]
\lean{Complexity.setupCfg_layout_return_step}
\label{Complexity.setupCfg_layout_return_step}
\leanok
\uses{Complexity.OblPhase, Complexity.OblSymbol, Complexity.nextThird, Complexity.obliviousSchedule, Complexity.setupCfg, Complexity.setupCfg_apply, Complexity.setupWrite, Turing.Action.apply, Turing.Cfg, Turing.Cfg.workTapeSymbols, Turing.FinTM, Turing.MultiTapeTM.step}
\difficulty{4}
Suppose a setup configuration is in layout-return phase $i \in \{0,1,2\}$ and the unary
tape does not read the origin sentinel at the unary head.  Then one step moves to
layout-return phase $i+1 \bmod 3$, moves the guide head one cell left, decrements the
unary head exactly when $i = 2$, and writes nothing on any tape.
\end{lemma}

\begin{lemma}[Prefix of the read-only return]
\lean{Complexity.setupCfg_layout_return_prefix}
\label{Complexity.setupCfg_layout_return_prefix}
\leanok
\uses{Complexity.OblPhase, Complexity.OblSymbol, Complexity.nextThird_mod, Complexity.obliviousSchedule, Complexity.setupCfg, Complexity.setupCfg_layout_return_step, Complexity.unaryTape, Complexity.unaryTape_origin, Turing.Cfg, Turing.FinTM, Turing.MultiTapeTM.runFrom, Turing.MultiTapeTM.runFrom_succ_eq_step'}
\difficulty{5}
Let $j \le 3B$.  From layout-return phase $0$ with unary head at cell $B$ over the
unary tape with $B$ markers and guide head at cell $3B$ (over an arbitrary guide tape),
exactly $j$ steps lead to layout-return phase $j \bmod 3$ with unary head at
$B - \lfloor j/3 \rfloor$ and guide head at $3B - j$; no tape is written and all guide
marks are unchanged.
\end{lemma}

\begin{lemma}[The completed return pass]
\lean{Complexity.setupCfg_layout_return}
\label{Complexity.setupCfg_layout_return}
\leanok
\uses{Complexity.OblPhase, Complexity.OblSymbol, Complexity.obliviousSchedule, Complexity.setupCfg, Complexity.setupCfg_apply, Complexity.setupCfg_layout_return_prefix, Complexity.setupWrite, Complexity.unaryTape, Turing.Action.apply, Turing.Cfg, Turing.Cfg.workTapeSymbols, Turing.FinTM, Turing.MultiTapeTM.runFrom, Turing.MultiTapeTM.runFrom_succ_eq_step'}
\difficulty{4}
From layout-return phase $0$ with unary head at cell $B$ over the unary tape with $B$
markers and guide head at cell $3B$ (over an arbitrary guide tape), exactly $3B+1$
steps reach the origin and enter the first leftward layout phase with unary head at
cell $1$ and guide head at cell $0$; no tape is written.
\end{lemma}

\begin{lemma}[The full rightward allocation pass]
\lean{Complexity.setupCfg_layout_right}
\label{Complexity.setupCfg_layout_right}
\leanok
\uses{Complexity.OblPhase, Complexity.OblSymbol, Complexity.copyGuide, Complexity.copyGuide_next, Complexity.copyGuide_origin, Complexity.guideFill_right_copy, Complexity.layoutMove, Complexity.layoutPhase, Complexity.obliviousSchedule, Complexity.rightGuide, Complexity.setupCfg, Complexity.setupCfg_layout_end, Complexity.setupCfg_layout_prefix, Complexity.setupCfg_right_edge, Complexity.unaryTape, Turing.Cfg, Turing.FinTM, Turing.MultiTapeTM.runFrom, Turing.MultiTapeTM.runFrom_succ_eq_step'}
\difficulty{5}
Let $n + 1 \le B$.  From the first rightward layout phase with unary head at cell $1$
over the unary tape with $B$ markers and guide head at cell $0$ over the partial guide
with $n+1$ copied cells, exactly $3B + 2$ steps reach radius $3B$, install the right
outer boundary marker, and enter the first layout-return phase with unary head at $B$,
guide head at $3B$, and guide tape the right guide of radius $3B$.
\end{lemma}

\begin{lemma}[The full leftward allocation pass]
\lean{Complexity.setupCfg_layout_left}
\label{Complexity.setupCfg_layout_left}
\leanok
\uses{Complexity.OblPhase, Complexity.OblSymbol, Complexity.guideFill, Complexity.guideFill_left_finish, Complexity.guideTape, Complexity.layoutMove, Complexity.layoutPhase, Complexity.obliviousSchedule, Complexity.rightGuide, Complexity.rightGuide_origin, Complexity.setupCfg, Complexity.setupCfg_apply, Complexity.setupCfg_layout_end, Complexity.setupCfg_layout_prefix, Complexity.setupWrite, Complexity.unaryTape, Turing.Action.apply, Turing.Cfg, Turing.FinTM, Turing.MultiTapeTM.runFrom, Turing.MultiTapeTM.runFrom_succ_eq_step', Turing.moveInputPos_zero}
\difficulty{5}
Let $B > 0$.  From the first leftward layout phase with unary head at cell $1$ over the
unary tape with $B$ markers and guide head at cell $0$ over the right guide of radius
$3B$, exactly $3B + 2$ steps complete the fixed guide and enter the unary-reset phase
with unary head at $B$, guide head at the left outer boundary $-3B - 1$, and guide tape
the completed guide of radius $3B$.
\end{lemma}

\begin{lemma}[Return scan to the guide origin]
\lean{Complexity.setupCfg_start_center}
\label{Complexity.setupCfg_start_center}
\leanok
\uses{Complexity.OblPhase, Complexity.OblSymbol, Complexity.guideTape, Complexity.guideTape_origin, Complexity.obliviousSchedule, Complexity.setupCfg, Complexity.setupCfg_apply, Complexity.setupWrite, Turing.Action.apply, Turing.Cfg, Turing.Cfg.workTapeSymbols, Turing.FinTM, Turing.MultiTapeTM.runFrom, Turing.MultiTapeTM.runFrom_succ_eq_step, Turing.MultiTapeTM.runFrom_zero, Turing.MultiTapeTM.step}
\difficulty{5}
For every $j \in \bbn$: from a setup configuration in the start-center phase with guide
head at cell $-j$ over the completed guide of radius $R$, exactly $j+1$ steps reach the
origin and enter the macrostep-check phase with guide head at cell $0$; all tapes and
all other heads are unchanged.
\end{lemma}

\begin{lemma}[Complete allocation ledger]
\lean{Complexity.setupCfg_allocation}
\label{Complexity.setupCfg_allocation}
\leanok
\uses{Complexity.OblPhase, Complexity.OblSymbol, Complexity.copyGuide, Complexity.guideTape, Complexity.obliviousSchedule, Complexity.setupCfg, Complexity.setupCfg_layout_left, Complexity.setupCfg_layout_return, Complexity.setupCfg_layout_right, Complexity.setupCfg_start_center, Complexity.setupCfg_unary_rewind, Complexity.setupCfg_unary_start, Complexity.unaryRewindPhase, Complexity.unaryTape, Turing.Cfg, Turing.FinTM, Turing.MultiTapeTM.runFrom, Turing.MultiTapeTM.runFrom_add}
\difficulty{5}
Let $n + 1 \le B$.  From a setup configuration in the unary-start phase with unary head
at cell $B+1$ over the unary tape with $B$ markers and guide head at cell $0$ over the
partial guide with $n+1$ copied cells, exactly $14B + 10$ steps produce the completed
guide of radius $3B$ and enter the macrostep-check phase with unary head at cell $1$
and guide head at cell $0$; the budget tape and head and the parked input head are
unchanged.
\end{lemma}

\begin{lemma}[The complete input-copy stage]
\lean{Complexity.setupCfg_copy}
\label{Complexity.setupCfg_copy}
\leanok
\uses{Complexity.OblPhase, Complexity.OblSymbol, Complexity.copyGuide, Complexity.obliviousSchedule, Complexity.setupCfg, Complexity.setupCfg_copy_end, Complexity.setupCfg_copy_left, Complexity.setupCfg_copy_prefix, Complexity.setupCfg_copy_return, Turing.Cfg, Turing.FinTM, Turing.MultiTapeTM.runFrom, Turing.MultiTapeTM.runFrom_add, Turing.MultiTapeTM.runFrom_succ_eq_step', Turing.MultiTapeTM.step}
\difficulty{4}
From a setup configuration in the copy-left phase with input head at position $1$,
guide head at cell $0$, and entirely blank guide tape, exactly $2|x| + 4$ steps ---
covering both virtual boundary cells and the return to the origin --- reach the
budget-start phase with input head parked at position $|x|+1$, guide head at cell $0$,
and guide tape the partial guide with $|x|+1$ copied cells; the budget and unary tapes
and heads are unchanged, and the native input stays parked from then on.
\end{lemma}

\begin{lemma}[Empty clock word equals empty unary budget]
\lean{Complexity.clockTape_unary_zero}
\label{Complexity.clockTape_unary_zero}
\leanok
\uses{Complexity.clockTape, Complexity.unaryTape, Turing.FinTM.bufferTape_nil}
\difficulty{2}
The captured clock tape of the empty bit word equals the unary tape with zero markers:
both consist of the origin sentinel at cell $0$ and blanks everywhere else.
\end{lemma}

\begin{lemma}[Setup form of a halted clock configuration]
\lean{Complexity.clockStageCfg_setup}
\label{Complexity.clockStageCfg_setup}
\leanok
\uses{Complexity.clockStageCfg, Complexity.clockTape, Complexity.clockTape_unary_zero, Complexity.finThree_cases, Complexity.setupCfg, Complexity.unaryTape, Turing.Cfg, Turing.FinTM}
\difficulty{4}
Let $c$ be a configuration of the $k$-tape Boolean machine $W$ on a Boolean input that
has halted (its state is the halting state).  Then the clock-stage embedding of $c$
into the oblivious schedule machine --- which keeps the work block of $c$ and captures
its entire output word on the budget tape --- is itself a setup configuration: base the
embedded configuration itself, phase reset-start, the embedded input-head position,
budget head at cell $m+1$ over the captured clock tape of the output word of $c$
(where $m$ is that word's length), unary head at cell $1$ over the unary tape with zero
markers, and guide head at cell $0$ over the blank tape.
\end{lemma}

\begin{lemma}[Macrostep retargeting of a setup configuration]
\lean{Complexity.macroCfg_setup}
\label{Complexity.macroCfg_setup}
\leanok
\uses{Complexity.OblPhase, Complexity.OblSymbol, Complexity.finThree_cases, Complexity.macroCfg, Complexity.setupCfg, Turing.Cfg, Turing.FinTM}
\difficulty{3}
Modifying a setup configuration by changing only its phase to $q'$ and its unary and
guide head positions to $u'$ and $g'$ --- keeping every tape, the budget head, and the
input head fixed --- again yields a setup configuration with the same remaining data.
\end{lemma}

\begin{lemma}[The completed setup runs out its unary budget and halts]
\lean{Complexity.setupCfg_finish}
\label{Complexity.setupCfg_finish}
\leanok
\uses{Complexity.OblPhase, Complexity.OblSymbol, Complexity.guideTape, Complexity.macroCfg_finish, Complexity.macroCfg_setup, Complexity.obliviousSchedule, Complexity.setupCfg, Complexity.unaryTape, Complexity.unaryTape_end, Complexity.unaryTape_unit, Turing.Cfg, Turing.FinTM, Turing.MultiTapeTM.runFrom}
\difficulty{3}
From a setup configuration in the macrostep-check phase with unary head at cell $1$
over the unary tape with $B$ markers and guide head at cell $0$ over the completed
guide of radius $3B$, exactly $B\,\bigl(6\,(3B) + 8\bigr) + 1$ steps execute one full
sweep per unary marker and reach the halted configuration with unary head at cell
$B+1$ and guide head at cell $0$; every tape, the budget head, and the input head are
unchanged.
\end{lemma}

\begin{lemma}[Complete initialization ledger]
\lean{Complexity.setupCfg_initializes}
\label{Complexity.setupCfg_initializes}
\leanok
\uses{Complexity.OblPhase, Complexity.OblSymbol, Complexity.budgetValue, Complexity.clockTape, Complexity.copyGuide, Complexity.guideTape, Complexity.obliviousSchedule, Complexity.setupCfg, Complexity.setupCfg_allocation, Complexity.setupCfg_budget_all, Complexity.setupCfg_copy, Complexity.setupCfg_reset, Complexity.unaryTape, Turing.Cfg, Turing.FinTM, Turing.MultiTapeTM.runFrom, Turing.MultiTapeTM.runFrom_add}
\difficulty{4}
Let $w$ be a bit word of little-endian value $v$, and suppose the input satisfies
$|x| + 1 \le (a+1)(v+1)$, where $a$ is the padding parameter.  Consider a setup
configuration in the reset-start phase with input head at an arbitrary position $p$,
budget head at cell $|w|+1$ over the captured clock tape of $w$, unary head at cell $1$
over the unary tape with zero markers, and guide head at cell $0$ over the blank guide
tape.  Then there is a bit word $w'$ with $|w'| = |w|$ such that after exactly
\[
  \bigl(\max(p,1)+1\bigr) + \bigl(2|x| + 4\bigr)
  + (v+1)\bigl(2|w| + (a+1) + 4\bigr) + \bigl(14\,(a+1)(v+1) + 10\bigr)
\]
steps the machine is in the macrostep-check phase with input head parked at position
$|x|+1$, budget head at cell $|w|+1$ over the captured clock tape of $w'$, unary head
at cell $1$ over the unary tape with $(a+1)(v+1)$ markers, and guide head at cell $0$
over the completed guide of radius $3\,(a+1)(v+1)$: the binary counter retains its full
width throughout.
\end{lemma}
